\documentclass[a4paper,oneside,11pt]{book}

\usepackage{graphicx}
\usepackage{titling}
\usepackage{geometry}
\usepackage{listings}
\usepackage{xcolor}
\usepackage{float}
\usepackage{hyperref}
\usepackage{indentfirst}

\geometry{
    left=3cm,
    right=2.5cm,
    top=3cm,
    bottom=3cm
}

\definecolor{backcolour}{rgb}{0.95,0.95,0.92}

\lstdefinestyle{phpstyle}{
    backgroundcolor=\color{backcolour},
    basicstyle=\ttfamily\footnotesize,
    breaklines=true,
    frame=single,
    language=PHP
}

\title{Laporan Praktikum \\ Praktikum Pemrograman Web 2\\
Pertemuan 9\\
Middleware Role-based Laravel}
\author{Nama Mahasiswa: Muhammad Rakan Hibrizi (24/539132/SV/24568)}
\date{\today}

\begin{document}

\begin{titlingpage}
\begin{center}
\vspace{4cm}
\Huge \textbf{\thetitle}\\[2cm]
\includegraphics[height=8cm]{lambang ugm.png}\\[4cm]
\Large
\theauthor\\
Kelas B2 \\[1cm]
Dosen Pengampu: Dinar Nugroho Pratomo, S.Kom., M.IM., M.Cs.\\[2cm]
\thedate
\end{center}
\end{titlingpage}

\tableofcontents
\newpage

\chapter{Tujuan Praktikum}
\begin{enumerate}
    \item Mahasiswa memahami konsep middleware Laravel dalam autentikasi dan otorisasi.
    \item Mahasiswa dapat membatasi akses halaman berdasarkan status login menggunakan middleware \texttt{auth}.
    \item Mahasiswa mampu membuat custom middleware untuk pengecekan role pengguna.
    \item Mahasiswa dapat mendaftarkan middleware custom di Laravel 11+.
    \item Mahasiswa mampu menerapkan middleware untuk membedakan hak akses antara user dan admin.
\end{enumerate}

\chapter{Tugas Praktikum}

\chapter{Tugas Praktikum}

\section{Langkah 1: Membuat Route Baru /hello}

Route sederhana untuk menampilkan pesan "Hello World".

\begin{lstlisting}[style=phpstyle]
// routes/web.php

Route::get('/hello', function () {
    return 'Hello World';
});
\end{lstlisting}

\textbf{Penjelasan:}
\begin{itemize}
    \item Route GET sederhana yang mengembalikan string
    \item Tidak memerlukan controller atau view
    \item Dapat diakses oleh siapa saja tanpa autentikasi
\end{itemize}

\section{Langkah 2: Membuat Controller dan Route}

\subsection{Membuat JobController}

Buat controller menggunakan artisan command:

\begin{lstlisting}[language=bash]
php artisan make:controller JobController
\end{lstlisting}

Tambahkan method di JobController:

\begin{lstlisting}[style=phpstyle]
// app/Http/Controllers/JobController.php

<?php

namespace App\Http\Controllers;

use Illuminate\Http\Request;

class JobController extends Controller
{
    public function index()
    {
        return 'Halaman Daftar Pekerjaan';
    }
}
\end{lstlisting}

\subsection{Mendaftarkan Route ke Controller}

\begin{lstlisting}[style=phpstyle]
// routes/web.php

use App\Http\Controllers\JobController;

Route::get('/jobs', [JobController::class, 'index']);
\end{lstlisting}

\textbf{Penjelasan:}
\begin{itemize}
    \item Controller memisahkan logic dari route
    \item Menggunakan array syntax \texttt{[ControllerClass::class, 'method']}
    \item Jangan lupa import class Controller di bagian atas file
\end{itemize}

\section{Langkah 3: Menggunakan Middleware Bawaan (auth)}

Middleware \texttt{auth} digunakan untuk membatasi akses hanya untuk user yang sudah login.

\begin{lstlisting}[style=phpstyle]
// routes/web.php

Route::get('/dashboard', function () {
    return view('dashboard');
})->middleware(['auth', 'verified'])->name('dashboard');
\end{lstlisting}

\textbf{Penjelasan:}
\begin{itemize}
    \item Middleware \texttt{auth} memastikan user sudah login
    \item Middleware \texttt{verified} memastikan email sudah diverifikasi
    \item Jika belum login, user akan di-redirect ke halaman login
    \item Jika belum verifikasi email, akan di-redirect ke halaman verifikasi
\end{itemize}

\section{Langkah 4: Membuat Custom Middleware IsAdmin}

\subsection{Generate Middleware}

Buat middleware menggunakan artisan command:

\begin{lstlisting}[language=bash]
php artisan make:middleware IsAdmin
\end{lstlisting}

\subsection{Implementasi Logic IsAdmin}

Edit file middleware yang telah dibuat:

\begin{lstlisting}[style=phpstyle]
// app/Http/Middleware/IsAdmin.php

<?php

namespace App\Http\Middleware;

use Closure;
use Illuminate\Http\Request;
use Symfony\Component\HttpFoundation\Response;
use Illuminate\Support\Facades\Auth;

class IsAdmin
{
    public function handle(Request $request, Closure $next): Response
    {
        // Check if user is authenticated and has admin role
        if (Auth::check() && Auth::user()->role === 'admin') {
            return $next($request);
        }

        // Redirect to home when not authorized
        return redirect('/')->with('error', 'Access denied. Admin only.');
    }
}
\end{lstlisting}

\textbf{Penjelasan:}
\begin{itemize}
    \item Method \texttt{handle()} menerima \texttt{\$request} dan \texttt{\$next}
    \item \texttt{Auth::check()} memeriksa apakah user sudah login
    \item \texttt{Auth::user()->role} mengambil role dari user yang login
    \item Jika role = 'admin', request diteruskan dengan \texttt{\$next(\$request)}
    \item Jika bukan admin, redirect ke home dengan pesan error
\end{itemize}

\section{Langkah 5: Mendaftarkan Middleware Custom}

Pada Laravel 11+, middleware didaftarkan di \texttt{bootstrap/app.php} menggunakan method \texttt{alias()}.

\begin{lstlisting}[style=phpstyle]
// bootstrap/app.php

<?php

use Illuminate\Foundation\Application;
use Illuminate\Foundation\Configuration\Exceptions;
use Illuminate\Foundation\Configuration\Middleware;

return Application::configure(basePath: dirname(__DIR__))
    ->withRouting(
        web: __DIR__.'/../routes/web.php',
        commands: __DIR__.'/../routes/console.php',
        health: '/up',
    )
    ->withMiddleware(function (Middleware $middleware): void {
        // Register middleware aliases
        $middleware->alias([
            'isAdmin' => \App\Http\Middleware\IsAdmin::class,
        ]);
    })
    ->withExceptions(function (Exceptions $exceptions): void {
        //
    })->create();
\end{lstlisting}

\textbf{Penjelasan:}
\begin{itemize}
    \item Middleware diberi alias 'isAdmin' agar mudah digunakan di route
    \item Method \texttt{->alias()} menerima array dengan format \texttt{'alias' => MiddlewareClass::class}
    \item Setelah didaftarkan, middleware dapat dipanggil dengan \texttt{->middleware('isAdmin')}
\end{itemize}

\textbf{Catatan:} Pada Laravel versi lama (sebelum Laravel 11), middleware didaftarkan di \texttt{app/Http/Kernel.php} pada array \texttt{\$routeMiddleware}.

\section{Langkah 6: Menambahkan Kolom Role ke Tabel Users}

\subsection{Modifikasi Migration}

Edit migration users table:

\begin{lstlisting}[style=phpstyle]
// database/migrations/0001_01_01_000000_create_users_table.php

Schema::create('users', function (Blueprint $table) {
    $table->id();
    $table->string('name');
    $table->string('email')->unique();
    $table->timestamp('email_verified_at')->nullable();
    $table->string('password');
    $table->enum('role', ['user', 'admin'])->default('user');
    $table->rememberToken();
    $table->timestamps();
});
\end{lstlisting}

\textbf{Penjelasan:}
\begin{itemize}
    \item Menambahkan kolom \texttt{role} dengan tipe ENUM
    \item Hanya menerima nilai 'user' atau 'admin'
    \item Default value adalah 'user'
\end{itemize}

\subsection{Jalankan Migration}

Karena migration sudah pernah dijalankan, gunakan \texttt{migrate:fresh} untuk reset database:

\begin{lstlisting}[language=bash]
php artisan migrate:fresh --seed
\end{lstlisting}

\textbf{Peringatan:} \texttt{migrate:fresh} akan menghapus semua data di database!

\subsection{Update Model User}

Tambahkan 'role' ke \texttt{\$fillable} agar dapat di-assign:

\begin{lstlisting}[style=phpstyle]
// app/Models/User.php

protected $fillable = [
    'name',
    'email',
    'password',
    'role',
];
\end{lstlisting}

\section{Langkah 7: Membatasi Route Khusus Admin}

Buat route yang hanya dapat diakses oleh admin dengan menggabungkan middleware \texttt{auth} dan \texttt{isAdmin}.

\begin{lstlisting}[style=phpstyle]
// routes/web.php

Route::get('/admin', function () {
    return 'Halaman Admin';
})->middleware(['auth', 'verified', 'isAdmin']);
\end{lstlisting}

\textbf{Penjelasan:}
\begin{itemize}
    \item Middleware \texttt{auth} memastikan user sudah login
    \item Middleware \texttt{verified} memastikan email sudah diverifikasi
    \item Middleware \texttt{isAdmin} memastikan role user adalah 'admin'
    \item Ketiga middleware dijalankan secara berurutan
    \item Jika salah satu middleware gagal, request tidak akan diteruskan
\end{itemize}

\chapter{Tugas Latihan}

\section{Latihan 1: Route /profile untuk User Login}

Buatlah route \texttt{/profile} yang hanya dapat diakses oleh user yang sudah login (tidak perlu admin).

\subsection{Solusi:}

\begin{lstlisting}[style=phpstyle]
// routes/web.php

Route::get('/profile', function () {
    return 'Halaman Profile: ' . Auth::user()->name;
})->middleware(['auth', 'verified'])->name('profile');
\end{lstlisting}

\textbf{Penjelasan:}
\begin{itemize}
    \item Menggunakan middleware \texttt{auth} dan \texttt{verified}
    \item Tidak menggunakan middleware \texttt{isAdmin}, sehingga semua user yang login dapat akses
    \item Menampilkan nama user dengan \texttt{Auth::user()->name}
\end{itemize}

\section{Latihan 2: Route /admin/jobs untuk Admin}

Buatlah route \texttt{/admin/jobs} yang hanya dapat diakses oleh admin.

\subsection{Solusi:}

\begin{lstlisting}[style=phpstyle]
// routes/web.php

Route::get('/admin/jobs', [AdminController::class, 'index'])
    ->middleware(['auth', 'verified', 'isAdmin'])
    ->name('admin.jobs');
\end{lstlisting}

Atau menggunakan route group untuk beberapa route admin:

\begin{lstlisting}[style=phpstyle]
// routes/web.php

Route::middleware(['auth', 'verified', 'isAdmin'])
    ->prefix('admin')
    ->group(function () {
        Route::get('/jobs', [AdminController::class, 'index'])
            ->name('admin.jobs');
        Route::get('/users', [AdminController::class, 'users'])
            ->name('admin.users');
        Route::get('/settings', [AdminController::class, 'settings'])
            ->name('admin.settings');
    });
\end{lstlisting}

\textbf{Penjelasan:}
\begin{itemize}
    \item Menggunakan middleware \texttt{auth}, \texttt{verified}, dan \texttt{isAdmin}
    \item \texttt{prefix('admin')} membuat semua route dalam group memiliki prefix '/admin'
    \item Semua route dalam group akan otomatis menggunakan ketiga middleware
    \item URL akan menjadi: /admin/jobs, /admin/users, /admin/settings
\end{itemize}

\chapter{Hasil dan Pembahasan}

\section{Testing Access Control}

\subsection{Test 1: Akses Route /hello}
\begin{itemize}
    \item URL: http://127.0.0.1:8000/hello
    \item Status: [\checkmark] Dapat diakses tanpa login
    \item Result: Menampilkan "Hello World"
\end{itemize}

\subsection{Test 2: Akses Route /jobs}
\begin{itemize}
    \item URL: http://127.0.0.1:8000/jobs
    \item Status: [\checkmark] Dapat diakses tanpa login
    \item Result: Menampilkan "Halaman Daftar Pekerjaan"
\end{itemize}

\subsection{Test 3: Akses Route /dashboard Tanpa Login}
\begin{itemize}
    \item URL: http://127.0.0.1:8000/dashboard
    \item Status: [$\times$] Redirect ke /login
    \item Result: Middleware \texttt{auth} berhasil memblokir akses
\end{itemize}

\subsection{Test 4: Akses Route /dashboard Setelah Login (User Biasa)}
\begin{itemize}
    \item Login sebagai: user@example.com
    \item URL: http://127.0.0.1:8000/dashboard
    \item Status: [\checkmark] Dapat diakses
    \item Result: Menampilkan dashboard
\end{itemize}

\subsection{Test 5: Akses Route /admin dengan User Biasa}
\begin{itemize}
    \item Login sebagai: user@example.com (role: user)
    \item URL: http://127.0.0.1:8000/admin
    \item Status: [$\times$] Redirect ke home dengan error message
    \item Result: Middleware \texttt{isAdmin} berhasil memblokir akses
\end{itemize}

\subsection{Test 6: Akses Route /admin dengan Admin}
\begin{itemize}
    \item Login sebagai: admin@example.com (role: admin)
    \item URL: http://127.0.0.1:8000/admin
    \item Status: [\checkmark] Dapat diakses
    \item Result: Menampilkan "Halaman Admin"
\end{itemize}

\subsection{Test 7: Akses Route /profile Setelah Login}
\begin{itemize}
    \item Login sebagai user atau admin
    \item URL: http://127.0.0.1:8000/profile
    \item Status: [\checkmark] Dapat diakses oleh user dan admin
    \item Result: Menampilkan "Halaman Profile: [Nama User]"
\end{itemize}

\subsection{Test 8: Akses Route /admin/jobs dengan Admin}
\begin{itemize}
    \item Login sebagai: admin@example.com
    \item URL: http://127.0.0.1:8000/admin/jobs
    \item Status: [\checkmark] Dapat diakses
    \item Result: Menampilkan halaman job management untuk admin
\end{itemize}

\section{Pembahasan}

\subsection{Cara Kerja Middleware}

Middleware bekerja sebagai filter yang memeriksa HTTP request sebelum mencapai route handler atau controller. Dalam praktikum ini:

\begin{enumerate}
    \item \textbf{Middleware auth}: Memeriksa apakah user sudah login. Jika belum, redirect ke halaman login.
    
    \item \textbf{Middleware verified}: Memeriksa apakah email user sudah diverifikasi. Jika belum, redirect ke halaman verifikasi email.
    
    \item \textbf{Middleware isAdmin}: Memeriksa role user. Jika bukan admin, redirect ke home dengan pesan error.
\end{enumerate}

\subsection{Urutan Eksekusi Middleware}

Ketika menggunakan multiple middleware seperti:

\begin{lstlisting}[style=phpstyle]
->middleware(['auth', 'verified', 'isAdmin'])
\end{lstlisting}

Urutan eksekusi adalah:
\begin{enumerate}
    \item Middleware \texttt{auth} dijalankan terlebih dahulu
    \item Jika \texttt{auth} pass, lanjut ke middleware \texttt{verified}
    \item Jika \texttt{verified} pass, lanjut ke middleware \texttt{isAdmin}
    \item Jika semua pass, request diteruskan ke controller/route handler
\end{enumerate}

Jika salah satu middleware gagal (tidak pass), eksekusi dihentikan dan request tidak diteruskan.

\subsection{Keuntungan Menggunakan Middleware}

\begin{enumerate}
    \item \textbf{Reusability}: Middleware dapat digunakan di banyak route tanpa perlu menulis ulang logic
    \item \textbf{Separation of Concerns}: Logic authorization terpisah dari controller
    \item \textbf{Maintainability}: Perubahan logic authorization hanya perlu dilakukan di satu tempat
    \item \textbf{Security}: Centralized security logic memudahkan audit
\end{enumerate}

\subsection{Perbedaan Laravel 11+ dengan Versi Sebelumnya}

\begin{table}[H]
\centering
\begin{tabular}{|l|l|l|}
\hline
\textbf{Aspek} & \textbf{Laravel 10 dan sebelumnya} & \textbf{Laravel 11+} \\
\hline
File registrasi & app/Http/Kernel.php & bootstrap/app.php \\
Method & \$routeMiddleware array & ->alias() method \\
Struktur & Array langsung & Closure function \\
\hline
\end{tabular}
\caption{Perbedaan Registrasi Middleware}
\end{table}

\chapter{Kesimpulan}

\subsection{Modifikasi Migration Users Table}
Menambahkan kolom \texttt{role} pada tabel users dengan tipe ENUM yang memiliki dua nilai: 'user' dan 'admin'.

\begin{lstlisting}[style=phpstyle]
// database/migrations/0001_01_01_000000_create_users_table.php

Schema::create('users', function (Blueprint $table) {
    $table->id();
    $table->string('name');
    $table->string('email')->unique();
    $table->timestamp('email_verified_at')->nullable();
    $table->string('password');
    $table->enum('role', ['user', 'admin'])->default('user');
    $table->rememberToken();
    $table->timestamps();
});
\end{lstlisting}

\subsection{Update UserFactory}
Menambahkan default role 'user' dan method \texttt{admin()} untuk membuat user dengan role admin.

\begin{lstlisting}[style=phpstyle]
// database/factories/UserFactory.php

public function definition(): array
{
    return [
        'name' => fake()->name(),
        'email' => fake()->unique()->safeEmail(),
        'email_verified_at' => now(),
        'password' => static::$password ??= Hash::make('password'),
        'role' => 'user',
        'remember_token' => Str::random(10),
    ];
}

public function admin(): static
{
    return $this->state(fn (array $attributes) => [
        'role' => 'admin',
    ]);
}
\end{lstlisting}

\subsection{Database Seeder}
Membuat test user dengan role admin dan user untuk keperluan testing.

\begin{lstlisting}[style=phpstyle]
// database/seeders/DatabaseSeeder.php

public function run(): void
{
    // Create regular users
    User::factory(5)->create();

    // Create admin user
    User::factory()->admin()->create([
        'name' => 'Admin User',
        'email' => 'admin@example.com',
        'password' => bcrypt('password'),
    ]);

    // Create regular test user
    User::factory()->create([
        'name' => 'Test User',
        'email' => 'user@example.com',
        'password' => bcrypt('password'),
        'role' => 'user',
    ]);
}
\end{lstlisting}

\section{Implementasi Custom Middleware}

\subsection{Middleware EnsureUserHasRole}
Middleware ini memeriksa apakah user memiliki role yang sesuai. Middleware ini dapat menerima multiple roles sebagai parameter.

\begin{lstlisting}[style=phpstyle]
// app/Http/Middleware/EnsureUserHasRole.php

<?php

namespace App\Http\Middleware;

use Closure;
use Illuminate\Http\Request;
use Illuminate\Support\Facades\Auth;
use Symfony\Component\HttpFoundation\Response;

class EnsureUserHasRole
{
    public function handle(Request $request, Closure $next, 
                          string ...$roles): Response
    {
        // Check if user is authenticated
        if (!Auth::check()) {
            return redirect()->route('login');
        }

        // Get user's role
        $userRole = Auth::user()->role;

        // Check if user's role matches any of the required roles
        if (in_array($userRole, $roles)) {
            return $next($request);
        }

        // If role doesn't match, abort with 403
        abort(403, 'Unauthorized access');
    }
}
\end{lstlisting}

\textbf{Penjelasan:}
\begin{itemize}
    \item Parameter \texttt{...\$roles} memungkinkan middleware menerima multiple roles
    \item Cek autentikasi terlebih dahulu menggunakan \texttt{Auth::check()}
    \item Bandingkan role user dengan array roles yang diizinkan menggunakan \texttt{in\_array()}
    \item Return 403 Forbidden jika role tidak sesuai
\end{itemize}

\subsection{Middleware IsAdmin}
Middleware khusus untuk memeriksa apakah user adalah admin. Jika bukan admin, user akan di-redirect ke halaman home.

\begin{lstlisting}[style=phpstyle]
// app/Http/Middleware/IsAdmin.php

<?php

namespace App\Http\Middleware;

use Closure;
use Illuminate\Http\Request;
use Symfony\Component\HttpFoundation\Response;
use Illuminate\Support\Facades\Auth;

class IsAdmin
{
    public function handle(Request $request, Closure $next): Response
    {
        // Check if user is authenticated and has admin role
        if (Auth::check() && Auth::user()->role === 'admin') {
            return $next($request);
        }

        // Redirect to home when not authorized
        return redirect('/')->with('error', 'Access denied. Admin only.');
    }
}
\end{lstlisting}

\textbf{Penjelasan:}
\begin{itemize}
    \item Khusus memeriksa role 'admin'
    \item Redirect ke home dengan flash message error jika bukan admin
    \item Lebih sederhana dibanding EnsureUserHasRole karena hanya check satu role
\end{itemize}

\subsection{Terminable Middleware - LogAfterRequest}
Middleware ini mencatat informasi request setelah response dikirim ke client, sehingga tidak memperlambat response time.

\begin{lstlisting}[style=phpstyle]
// app/Http/Middleware/LogAfterRequest.php

<?php

namespace App\Http\Middleware;

use Closure;
use Illuminate\Http\Request;
use Illuminate\Support\Facades\Log;
use Symfony\Component\HttpFoundation\Response;

class LogAfterRequest
{
    public function handle(Request $request, Closure $next): Response
    {
        return $next($request);
    }

    public function terminate(Request $request, Response $response): void
    {
        // Log request details after response is sent
        Log::info('Request processed', [
            'url' => $request->fullUrl(),
            'method' => $request->method(),
            'ip' => $request->ip(),
            'user_id' => auth()->id(),
            'status' => $response->getStatusCode(),
            'timestamp' => now()->toDateTimeString(),
        ]);
    }
}
\end{lstlisting}

\textbf{Penjelasan:}
\begin{itemize}
    \item Method \texttt{terminate()} dipanggil setelah response dikirim
    \item Logging tidak mempengaruhi response time
    \item Mencatat URL, method HTTP, IP address, user ID, status code, dan timestamp
    \item Berguna untuk audit trail dan debugging
\end{itemize}

\section{Registrasi Middleware}

Pada Laravel 11+, middleware didaftarkan menggunakan method \texttt{alias()} di file \texttt{bootstrap/app.php}.

\begin{lstlisting}[style=phpstyle]
// bootstrap/app.php

<?php

use Illuminate\Foundation\Application;
use Illuminate\Foundation\Configuration\Exceptions;
use Illuminate\Foundation\Configuration\Middleware;

return Application::configure(basePath: dirname(__DIR__))
    ->withRouting(
        web: __DIR__.'/../routes/web.php',
        commands: __DIR__.'/../routes/console.php',
        health: '/up',
    )
    ->withMiddleware(function (Middleware $middleware): void {
        // Register middleware aliases
        $middleware->alias([
            'role' => \App\Http\Middleware\EnsureUserHasRole::class,
            'isAdmin' => \App\Http\Middleware\IsAdmin::class,
            'log.after' => \App\Http\Middleware\LogAfterRequest::class,
        ]);
    })
    ->withExceptions(function (Exceptions $exceptions): void {
        //
    })->create();
\end{lstlisting}

\textbf{Penjelasan:}
\begin{itemize}
    \item Setiap middleware diberi alias untuk memudahkan penggunaan di routes
    \item Alias 'role' untuk EnsureUserHasRole
    \item Alias 'isAdmin' untuk IsAdmin
    \item Alias 'log.after' untuk LogAfterRequest
\end{itemize}

\section{Implementasi Routes dengan Middleware}

\subsection{Route Structure}
Routes diorganisir berdasarkan level akses dan role pengguna.

\begin{lstlisting}[style=phpstyle]
// routes/web.php

<?php

use App\Http\Controllers\ProfileController;
use App\Http\Controllers\AdminController;
use Illuminate\Support\Facades\Route;

// Public route
Route::get('/', function () {
    return view('welcome');
});

// Dashboard untuk semua user yang login
Route::get('/dashboard', function () {
    return view('dashboard');
})->middleware(['auth', 'verified'])->name('dashboard');
\end{lstlisting}

\subsection{User Routes}
Routes khusus untuk user dengan role 'user'. Menggunakan middleware \texttt{role:user} untuk membatasi akses.

\begin{lstlisting}[style=phpstyle]
// User Dashboard - hanya untuk role 'user'
Route::middleware(['auth', 'verified', 'role:user'])
    ->prefix('user')
    ->group(function () {
        Route::get('/dashboard', function () {
            return view('user.dashboard');
        })->name('user.dashboard');
        
        Route::get('/profile', function () {
            return view('user.profile');
        })->name('user.profile');
    });
\end{lstlisting}

\textbf{Penjelasan:}
\begin{itemize}
    \item \texttt{middleware(['auth', 'verified', 'role:user'])} - Harus login, email verified, dan role user
    \item \texttt{prefix('user')} - Semua route dalam group ini akan memiliki prefix '/user'
    \item Hanya user dengan role 'user' yang dapat mengakses
\end{itemize}

\subsection{Admin Routes}
Routes khusus untuk admin menggunakan middleware \texttt{isAdmin} dan \texttt{log.after} untuk logging.

\begin{lstlisting}[style=phpstyle]
// Admin Dashboard - hanya untuk role 'admin'
Route::middleware(['auth', 'verified', 'isAdmin', 'log.after'])
    ->prefix('admin')
    ->group(function () {
        Route::get('/jobs', [AdminController::class, 'index'])
            ->name('admin.jobs');
        Route::get('/users', [AdminController::class, 'users'])
            ->name('admin.users');
        Route::get('/settings', [AdminController::class, 'settings'])
            ->name('admin.settings');
        Route::get('/reports', [AdminController::class, 'reports'])
            ->name('admin.reports');
    });
\end{lstlisting}

\textbf{Penjelasan:}
\begin{itemize}
    \item \texttt{isAdmin} middleware memastikan hanya admin yang dapat akses
    \item \texttt{log.after} mencatat setiap request admin untuk audit
    \item Semua route menggunakan AdminController
    \item URL akan menjadi /admin/jobs, /admin/users, dll
\end{itemize}

\section{Implementasi Controller}

\subsection{AdminController}
Controller untuk menangani semua request admin.

\begin{lstlisting}[style=phpstyle]
// app/Http/Controllers/AdminController.php

<?php

namespace App\Http\Controllers;

use Illuminate\Http\Request;
use App\Models\User;

class AdminController extends Controller
{
    public function index()
    {
        return view('admin.jobs', [
            'title' => 'Admin - Job Management'
        ]);
    }

    public function users()
    {
        $users = User::all();
        return view('admin.users', [
            'users' => $users,
            'title' => 'User Management'
        ]);
    }

    public function settings()
    {
        return view('admin.settings', [
            'title' => 'Admin Settings'
        ]);
    }

    public function reports()
    {
        return view('admin.reports', [
            'title' => 'Reports'
        ]);
    }
}
\end{lstlisting}

\subsection{Update RegisteredUserController}
Menambahkan validasi dan penyimpanan role saat registrasi.

\begin{lstlisting}[style=phpstyle]
// app/Http/Controllers/Auth/RegisteredUserController.php

public function store(Request $request): RedirectResponse
{
    $request->validate([
        'name' => ['required', 'string', 'max:255'],
        'email' => ['required', 'string', 'lowercase', 'email', 
                   'max:255', 'unique:'.User::class],
        'password' => ['required', 'confirmed', 
                      Rules\Password::defaults()],
        'role' => ['required', 'in:user,admin'],
    ]);

    $user = User::create([
        'name' => $request->name,
        'email' => $request->email,
        'password' => Hash::make($request->password),
        'role' => $request->role,
    ]);

    event(new Registered($user));
    Auth::login($user);

    return redirect(route('dashboard', absolute: false));
}
\end{lstlisting}

\textbf{Penjelasan:}
\begin{itemize}
    \item Validasi \texttt{in:user,admin} memastikan role hanya 'user' atau 'admin'
    \item Role disimpan saat user register
    \item Setelah register, user langsung login dan redirect ke dashboard
\end{itemize}

\section{Implementasi Views}

\subsection{User Dashboard}
View untuk user dengan role 'user'.

\begin{lstlisting}[style=phpstyle]
<!-- resources/views/user/dashboard.blade.php -->

<x-app-layout>
    <x-slot name="header">
        <h2 class="font-semibold text-xl text-gray-800 leading-tight">
            {{ __('User Dashboard') }}
        </h2>
    </x-slot>

    <div class="py-12">
        <div class="max-w-7xl mx-auto sm:px-6 lg:px-8">
            <div class="bg-white overflow-hidden shadow-sm sm:rounded-lg">
                <div class="p-6 text-gray-900">
                    <h3 class="text-2xl font-bold mb-4">
                        Welcome, {{ Auth::user()->name }}!
                    </h3>
                    <p class="text-gray-600 mb-6">
                        You are logged in as: 
                        <span class="font-semibold text-blue-600">
                            {{ Auth::user()->role }}
                        </span>
                    </p>
                    
                    <!-- User content here -->
                </div>
            </div>
        </div>
    </div>
</x-app-layout>
\end{lstlisting}

\subsection{Admin Jobs View}
View untuk admin job management.

\begin{lstlisting}[style=phpstyle]
<!-- resources/views/admin/jobs.blade.php -->

<x-app-layout>
    <x-slot name="header">
        <h2 class="font-semibold text-xl text-gray-800 leading-tight">
            {{ __('Admin - Job Management') }}
        </h2>
    </x-slot>

    <div class="py-12">
        <div class="max-w-7xl mx-auto sm:px-6 lg:px-8">
            <div class="bg-white overflow-hidden shadow-sm sm:rounded-lg">
                <div class="p-6 text-gray-900">
                    <h3 class="text-2xl font-bold mb-4">
                        Welcome, Admin {{ Auth::user()->name }}!
                    </h3>
                    
                    <!-- Admin content here -->
                </div>
            </div>
        </div>
    </div>
</x-app-layout>
\end{lstlisting}

\section{Testing Access Control}

\subsection{Test Scenario 1: Login sebagai User}
\begin{itemize}
    \item \textbf{Credentials:} user@example.com / password
    \item \textbf{Role:} user
    \item \textbf{Result:}
    \begin{itemize}
        \item [\checkmark] Dapat akses /user/dashboard
        \item [\checkmark] Dapat akses /user/profile
        \item [$\times$] Tidak dapat akses /admin/jobs (403 Forbidden)
        \item [$\times$] Tidak dapat akses /admin/users (redirect ke home)
    \end{itemize}
\end{itemize}

\subsection{Test Scenario 2: Login sebagai Admin}
\begin{itemize}
    \item \textbf{Credentials:} admin@example.com / password
    \item \textbf{Role:} admin
    \item \textbf{Result:}
    \begin{itemize}
        \item [\checkmark] Dapat akses /admin/jobs
        \item [\checkmark] Dapat akses /admin/users
        \item [\checkmark] Dapat akses /admin/settings
        \item [\checkmark] Dapat akses /admin/reports
        \item [$\times$] Tidak dapat akses /user/dashboard (403 Forbidden)
    \end{itemize}
\end{itemize}

\subsection{Test Scenario 3: Access tanpa Login}
\begin{itemize}
    \item \textbf{Result:}
    \begin{itemize}
        \item [$\times$] Semua protected routes redirect ke /login
        \item [$\times$] Session menyimpan intended URL
        \item [\checkmark] Setelah login, redirect ke intended URL
    \end{itemize}
\end{itemize}

\section{Log Output}

Contoh log yang dihasilkan oleh LogAfterRequest middleware:

\begin{lstlisting}[style=phpstyle]
// storage/logs/laravel.log

[2025-10-28 10:15:23] local.INFO: Request processed 
{
    "url": "http://127.0.0.1:8000/admin/jobs",
    "method": "GET",
    "ip": "127.0.0.1",
    "user_id": 9,
    "status": 200,
    "timestamp": "2025-10-28 10:15:23"
}
\end{lstlisting}

\section{Pembahasan}

\subsection{Keuntungan Middleware}

\subsubsection{Separation of Concerns}
Middleware memisahkan logic authorization dari controller, membuat code lebih maintainable dan mengikuti prinsip Single Responsibility Principle (SRP).

\subsubsection{Reusability}
Middleware dapat digunakan di banyak route tanpa perlu menulis ulang code. Cukup menambahkan alias middleware di route definition.

\subsubsection{Security}
Centralized security logic memudahkan audit dan maintenance. Semua logic authorization ada di satu tempat.

\subsubsection{Performance}
Terminable middleware tidak memperlambat response karena dieksekusi setelah response dikirim ke client.

\subsection{Perbedaan Middleware dengan Gate/Policy}

\begin{table}[H]
\centering
\begin{tabular}{|l|l|l|}
\hline
\textbf{Aspek} & \textbf{Middleware} & \textbf{Gate/Policy} \\
\hline
Level & HTTP Request & Authorization Logic \\
Timing & Before/After Request & On-demand \\
Use Case & Route protection & Resource authorization \\
Complexity & Simple checks & Complex business logic \\
\hline
\end{tabular}
\caption{Perbandingan Middleware dan Gate/Policy}
\end{table}

\subsection{Best Practices}

\subsubsection{Kapan Menggunakan Middleware}
\begin{itemize}
    \item Authentication checks
    \item Role-based access control
    \item Request logging dan monitoring
    \item Rate limiting
    \item CORS handling
\end{itemize}

\subsubsection{Kapan Menggunakan Gate/Policy}
\begin{itemize}
    \item Resource ownership checks (user dapat edit post miliknya sendiri)
    \item Complex authorization logic dengan multiple conditions
    \item Model-specific permissions
\end{itemize}

\subsubsection{Naming Convention}
\begin{itemize}
    \item Middleware class: PascalCase dengan prefix yang jelas (EnsureUserHasRole, IsAdmin)
    \item Middleware alias: lowercase dengan dot notation untuk grouping (role, isAdmin, log.after)
\end{itemize}

\section{Hasil Tampilan}

\begin{figure}[H]
    \centering
    \includegraphics[width=0.8\linewidth]{403_forbidden.png}
    \caption{Akses Ditolak dengan 403 Forbidden jika Role Tidak Sesuai}
\end{figure}

\begin{figure}[H]
    \centering
    \includegraphics[width=0.8\linewidth]{user_dashboard.png}
    \caption{User Dashboard - Hanya Dapat Diakses oleh Role User}
\end{figure}

\begin{figure}[H]
    \centering
    \includegraphics[width=0.8\linewidth]{admin_jobs.png}
    \caption{Admin Jobs - Hanya Dapat Diakses oleh Role Admin}
\end{figure}

\begin{figure}[H]
    \centering
    \includegraphics[width=0.8\linewidth]{register_form.png}
    \caption{Form Registrasi dengan Pilihan Role}
\end{figure}

\chapter{Kesimpulan}

\begin{enumerate}
    \item Middleware adalah filter yang bekerja di antara HTTP request dan response, memungkinkan kita untuk memeriksa dan memfilter request sebelum mencapai controller.
    
    \item Middleware \texttt{auth} bawaan Laravel sangat berguna untuk membatasi akses halaman hanya untuk user yang sudah login.
    
    \item Custom middleware seperti \texttt{IsAdmin} dapat dibuat untuk kebutuhan spesifik seperti pengecekan role atau permissions.
    
    \item Pada Laravel 11+, middleware didaftarkan di \texttt{bootstrap/app.php} menggunakan method \texttt{->alias()}, berbeda dengan Laravel 10 yang menggunakan \texttt{app/Http/Kernel.php}.
    
    \item Multiple middleware dapat digabungkan pada satu route, dan akan dieksekusi secara berurutan sesuai urutan pendaftaran.
    
    \item Route grouping dengan \texttt{prefix()} dan \texttt{group()} memudahkan penerapan middleware yang sama pada banyak route sekaligus.
\end{enumerate}

\chapter{Saran}

\begin{enumerate}
    \item Tambahkan lebih banyak role seperti 'moderator', 'editor' untuk sistem yang lebih kompleks.
    
    \item Buat middleware yang lebih fleksibel dengan parameter untuk mengecek multiple roles.
    
    \item Implementasikan unit testing untuk middleware menggunakan PHPUnit atau Pest.
    
    \item Gunakan package seperti Spatie Permission untuk role dan permission management yang lebih advanced.
    
    \item Tambahkan logging untuk tracking siapa saja yang mencoba mengakses route admin.
\end{enumerate}

\chapter{Daftar Pustaka}

\begin{enumerate}
    \item Laravel Documentation. (2024). \textit{Middleware}. Retrieved from https://laravel.com/docs/11.x/middleware
    
    \item Laravel Documentation. (2024). \textit{Authentication}. Retrieved from https://laravel.com/docs/11.x/authentication
    
    \item Laravel Documentation. (2024). \textit{Routing}. Retrieved from https://laravel.com/docs/11.x/routing
\end{enumerate}

\chapter{Lampiran}

\section{Daftar Route Lengkap}

\begin{table}[H]
\centering
\begin{tabular}{|l|l|l|l|}
\hline
\textbf{Method} & \textbf{URI} & \textbf{Middleware} & \textbf{Akses} \\
\hline
GET & /hello & - & Public \\
GET & /jobs & - & Public \\
GET & /dashboard & auth, verified & Logged in \\
GET & /admin & auth, verified, isAdmin & Admin only \\
GET & /profile & auth, verified & Logged in \\
GET & /admin/jobs & auth, verified, isAdmin & Admin only \\
\hline
\end{tabular}
\caption{Daftar Routes dengan Middleware}
\end{table}

\section{Test Credentials}

\begin{table}[H]
\centering
\begin{tabular}{|l|l|l|}
\hline
\textbf{Role} & \textbf{Email} & \textbf{Password} \\
\hline
Admin & admin@example.com & password \\
User & user@example.com & password \\
\hline
\end{tabular}
\caption{Test Credentials untuk Testing}
\end{table}

\section{Command Artisan yang Digunakan}

\begin{lstlisting}[language=bash]
# Membuat middleware
php artisan make:middleware IsAdmin

# Membuat controller
php artisan make:controller JobController

# Reset database dan seed
php artisan migrate:fresh --seed

# Melihat daftar routes
php artisan route:list

# Menjalankan server development
php artisan serve
\end{lstlisting}

\chapter{Link Github}

Repository: \url{https://github.com/RakanAja05/job-portal}

Branch: Middleware

Link Branch: \url{https://github.com/RakanAja05/job-portal/tree/Middleware}

\end{document}
